\section{双等腰模型的三项结论}
\block{共用构型}
如图，\(\triangle CBC'\) 和 \(\triangle DBD'\) 都是以 \(B\) 为直角顶点的等腰直角三角形，即
\[
BC=BC',\quad BD=BD',\quad
\angle CBC'=\angle DBD'=90^\circ.
\]
各点方向按图取定：射线 \(BD\)、\(BD'\)、\(BC\)、\(BC'\) 绕 \(B\) 依次排列，\(\angle CBD'\) 为锐角。从 \(BC\) 到 \(BC'\)、从 \(BD\) 到 \(BD'\) 的 \(90^\circ\) 转动方向相同。以下三题共用这一构型，分别作辅助线。
\diagram{\diModel}

\begin{problem}{题 1 等面积从哪里来}
求证：\(S_{\triangle CBD'}=S_{\triangle C'BD}\)。
\hint{已经有 \(BC=BC'\)。分别以 \(BC\)、\(BC'\) 为底，试着比较两条高。}
\end{problem}

\begin{problem}{题 2 已知中点，证明垂直}
\(E\) 是 \(CD'\) 的中点。连接 \(BE\)，求证：\(BE\perp C'D\)。
\hint{中点条件能否通过倍长 \(BE\)，转化成一组八字形全等？}
\end{problem}

\begin{problem}{题 3 已知垂直，证明中点}
过 \(B\) 作 \(BF\perp C'D\)，\(F\) 为垂足；直线 \(BF\) 与 \(CD'\) 相交于 \(E\)。求证：\(CE=ED'\)。
\hint{还不知道 \(E\) 是中点。分别从 \(C\)、\(D'\) 向直线 \(BF\) 作垂线，看看会出现几组一线三等角。}
\end{problem}
\clearpage
